\documentclass[11pt,letterpaper]{article}
\usepackage[top=0.7in,bottom=0.7in,left=0.75in,right=0.75in]{geometry}
\usepackage[T1]{fontenc}
\usepackage[utf8]{inputenc}
\usepackage{lmodern}
\usepackage{microtype}
\usepackage{enumitem}
\usepackage[hidelinks]{hyperref}
\usepackage{titlesec}
\setlength{\parindent}{0pt}
\setlength{\parskip}{0pt}
\pagestyle{empty}
\titleformat{\section}{\normalfont\bfseries}{}{0em}{\uppercase}[\vspace{-4pt}\hrule\vspace{4pt}]
\titlespacing{\section}{0pt}{8pt}{5pt}
\setlist[itemize]{leftmargin=1.2em,label=--,itemsep=0pt,topsep=2pt,parsep=0pt}
\begin{document}
\begin{center}
{\fontsize{16}{19}\selectfont\textbf{\MakeUppercase{Diego Castillo Pineda}}}\\[4pt]
{\small La Florida, Santiago, Chile \quad|\quad +56 9 94162547 \quad|\quad \href{mailto:diego.castillop11@gmail.com}{diego.castillop11@gmail.com}\\
\href{https://linkedin.com/in/diego-castillo-pineda}{linkedin.com/in/diego-castillo-pineda} \quad|\quad \href{https://github.com/diegocastillop}{github.com/diegocastillop}}
\end{center}

\section{Resumen Profesional}
Analista de Datos con sólida trayectoria en SQL Server, Azure y visualización de datos mediante Power BI, enfocado en convertir datos complejos en insights accionables para la toma de decisiones estratégicas. Experiencia directa en análisis de bases de datos de alta concurrencia, automatización de procesos mediante Python y T-SQL, y desarrollo de dashboards interactivos que optimizan el seguimiento de KPIs operativos. En Punto Ticket, diseñé procedimientos almacenados que procesaron más de 50,000 transacciones por hora durante eventos masivos, reduciendo tiempos de respuesta en 40\% y mejorando la estabilidad operativa en períodos críticos.

\section{Experiencia Profesional}

\textbf{Punto Ticket} \hfill Santiago, Chile\\
\textit{Analista de Base de Datos Senior} \hfill Nov. 2019 – Feb. 2026
\begin{itemize}
  \item Diseñé y optimicé stored procedures en SQL Server para eventos masivos, procesando más de 50,000 transacciones por hora y reduciendo tiempos de respuesta en 40\% durante períodos de alta demanda.
  \item Automaticé procesos de carga y transformación de datos mediante scripts T-SQL y Python, eliminando errores manuales en operaciones recurrentes y mejorando la eficiencia operativa en 35\%.
  \item Desarrollé consultas dinámicas con PIVOT y JSON para reportería avanzada, reemplazando procesos manuales y reduciendo el tiempo de generación de reportes de 4 horas a 15 minutos.
  \item Participé activamente en la migración de bases de datos hacia Microsoft Azure SQL, asegurando disponibilidad continua y mejorando la escalabilidad de la infraestructura cloud.
  \item Integré datos entre múltiples sistemas utilizando JSON y XML, facilitando la consolidación de información y el análisis de datos para decisiones de negocio.
  \item Administré bases de datos SQL Server en entornos de alta concurrencia, implementando monitoreo proactivo y planes de contingencia que garantizaron 99.8\% de disponibilidad en eventos críticos.
\end{itemize}
\vspace{2pt}
\textbf{Imperial S.A.} \hfill Santiago, Chile\\
\textit{Analista Funcional} \hfill Jun. 2017 – Nov. 2019
\begin{itemize}
  \item Realicé análisis funcional de sistemas ERP (Oracle y SAP), levantando requerimientos con usuarios clave y traduciendo necesidades de negocio en especificaciones técnicas para el equipo de desarrollo.
  \item Ejecuté consultas SQL complejas en Oracle y SQL Server para extracción, análisis y validación de datos, generando reportes personalizados que apoyaron la toma de decisiones operativas y financieras.
  \item Documenté procesos técnicos mediante manuales de usuario, diagramas de flujo e instructivos, facilitando la capacitación de nuevos usuarios y mejorando la eficiencia del soporte a 120+ usuarios finales.
  \item Proporcioné soporte técnico-funcional en ambientes productivos, resolviendo incidencias críticas en tiempo real y reduciendo el tiempo promedio de resolución en 25\%.
  \item Parametricé y configuré módulos de sistemas según especificaciones funcionales, participando en pruebas de calidad y validación de resultados antes de implementaciones productivas.
\end{itemize}
\vspace{2pt}
\textbf{OB GROUP PARK} \hfill Santiago, Chile\\
\textit{Analista de Sistemas y TI} \hfill Mayo 2014 – Abril 2017
\begin{itemize}
  \item Lideré proyectos de integración tecnológica para optimizar procesos contables y de gestión en el sector retail, coordinando la implementación de nuevas plataformas y garantizando compatibilidad entre sistemas heredados.
  \item Colaboré con áreas de negocio para identificar oportunidades de mejora mediante análisis de datos operativos, proponiendo soluciones tecnológicas que aumentaron la eficiencia de procesos administrativos en 20\%.
  \item Administré bases de datos SQL Server y proporcioné soporte técnico a usuarios finales, asegurando la continuidad operativa y la correcta adopción de cambios tecnológicos en múltiples sucursales.
\end{itemize}
\vspace{2pt}
\textbf{Hewlett-Packard Company} \hfill Santiago, Chile\\
\textit{Agente de Mesa de Ayuda Nivel 1} \hfill Ago. 2013 – Mayo 2014
\begin{itemize}
  \item Proporcioné soporte técnico de primer nivel a usuarios finales, diagnosticando y resolviendo incidencias de hardware, software y conectividad, desarrollando habilidades fundamentales en resolución de problemas y comunicación técnica efectiva.
\end{itemize}

\section{Proyectos Destacados}

\textbf{Sistema de Análisis Predictivo de Salud con IA} \hfill 2024
\begin{itemize}
  \item Desarrollé aplicación web con React y Node.js que utiliza Gemini AI para analizar patrones de datos médicos y generar insights predictivos mediante prompt engineering avanzado.
  \item Implementé base de datos MongoDB con encriptación de datos sensibles, garantizando cumplimiento de normativas de privacidad y seguridad en el manejo de información de pacientes.
  \item Diseñé UI/UX intuitiva con dashboards interactivos en React que visualizan métricas de salud en tiempo real, mejorando la experiencia de usuario y facilitando la interpretación de datos complejos.
  \item Integré pipeline de datos automatizado con Python para limpieza, transformación y carga (ETL) de registros médicos, procesando más de 10,000 registros diarios con precisión del 98\%.
\end{itemize}

\section{Habilidades T\'{e}cnicas}
\textbf{Análisis de Datos \& BI:} SQL Server (Experto), Azure SQL, Power BI, Excel Avanzado, DAX, Business Analytics, Data Visualization \\
\textbf{Bases de Datos \& Cloud:} SQL Server Management Studio, Oracle, PostgreSQL, MongoDB, Microsoft Azure, Azure Data Factory \\
\textbf{Desarrollo \& Automatización:} Python, T-SQL, Stored Procedures, ETL, C\#, Node.js, Git (GitHub/GitLab), Docker \\
\textbf{IA \& Desarrollo Frontend:} Prompt Engineering (Claude, Gemini), React, TypeScript, JavaScript (ES6+), HTML5/CSS3, JSON/XML, UI/UX Design

\section{Formaci\'{o}n Acad\'{e}mica}
\textbf{Business Analytics Essentials \& Data Science}, Universidad de Chile, Escuela de Postgrado \hfill 2023 - 2024 \\
\textbf{Administración de Bases de Datos SQL Server}, Universidad de Chile, Escuela de Postgrado \hfill 2023 \\
\textbf{Desarrollo de Apps Móviles}, Universidad Complutense de Madrid \hfill 2017 \\
\textbf{Analista Programador}, Universidad Tecnológica de Chile INACAP \hfill 2013

\end{document}